\documentclass[11pt]{article}
\usepackage[utf8]{inputenc}
\usepackage[T1]{fontenc}
\usepackage[a4paper, margin=1in]{geometry}
\usepackage{setspace}
\usepackage{microtype}
\usepackage{lmodern}

\pagestyle{empty}
\setlength{\parindent}{0pt}

\begin{document}

\begin{center}
{\Large\bfseries Amending Amenables: Labor Law and the\\[2pt]
Direction of Technical Change\par}

\vspace{14pt}

Kabeer Bora\\[2pt]
\textit{University of Missouri--Kansas City}

\vspace{18pt}

\textbf{ABSTRACT}
\end{center}

\vspace{4pt}

India's 1982 amendment to the Industrial Disputes Act required firms with 100 or more workers to obtain permission for layoffs and closure, making labor quasi-fixed. Using Annual Survey of Industries data and a difference-in-discontinuities design at the 100-worker threshold, we find treated private firms raised capital intensity by 15\% and cut capital productivity by 14\% in 1983, with labor productivity unchanged through 1986. The real wage did not move; the cost of adjusting labor did. This is defensive capital deepening: capital-using restructuring without the labor-saving gain of canonical Marx-biased technical change.

\end{document}
